\section{Introduction}
Higher education institutions (HEIs) have been of increasing importance in recent years, given their role in developing a more sustainable society. Therefore, preventing dropout and ensuring students' receive proper support is crucial. In this paper, we propose an approach based on decision tree (DT) using students' examination results to allocate assistance effectively.

Using multiple DT algorithms, we analyze a given dataset of students' examination results to identify patterns and factors that contribute to academic success or failure in order to provide an appropriate level of academic support. The DT-based approach allows us to classify students into different categories based on their performance, enabling targeted interventions and resource allocation. By identifying at-risk students early on, HEIs can implement strategies to improve retention rates and enhance the overall educational experience. The results of this study can inform policy decisions and contribute to the development of more effective support systems for students in HEIs.

Section II of this paper will attempt to explore key technical aspects of this problem, including different types of DT algorithms used, evaluation metrics, and techniques for handling imbalanced datasets. Section III will outline related work in the field of educational data mining and DT-based prediction models. Section IV and V will describe present the methodology and experimental setup, including data preprocessing, feature selection, and model training. Finally, Section VI will discuss the results and implications of our findings, as well as potential future research directions in this area.